\documentclass[11pt,a4paper,oneside,openany]{ltjsbook}

% --- Typography & Geometry ---
\usepackage{luatexja}
\usepackage{luatexja-fontspec}
\usepackage{amsmath,amssymb,amsthm}
\usepackage{geometry}
\geometry{top=24mm,bottom=24mm,left=24mm,right=24mm,footskip=10mm}

% --- Visual & Layout ---
\usepackage{tcolorbox}
\tcbuselibrary{skins,breakable}
\usepackage{tikz}
\usetikzlibrary{arrows.meta,calc,positioning,shapes.geometric}
\usepackage{graphicx}
\usepackage{xcolor}
\usepackage{booktabs}
\usepackage{fancyhdr}

% --- Color Palette ---
\definecolor{primaryblue}{RGB}{20, 50, 110}     % 深遠な紺青
\definecolor{accentcrimson}{RGB}{165, 35, 45}   % 格調ある深紅
\definecolor{softblue}{RGB}{240, 245, 252}      % 淡い背景青
\definecolor{darkslate}{RGB}{45, 55, 72}        % 濃墨
\definecolor{forestgreen}{RGB}{25, 115, 65}     % 深緑
\definecolor{warmgold}{RGB}{195, 140, 20}       % 琥珀金

% --- Puzzle Piece Colors for Haberdasher Dissection ---
\definecolor{pieceA}{RGB}{218, 232, 252}        % 澄んだスカイブルー
\definecolor{pieceB}{RGB}{255, 242, 204}        % 温もりある琥珀イエロー
\definecolor{pieceC}{RGB}{213, 232, 212}        % 爽やかなミントグリーン
\definecolor{pieceD}{RGB}{248, 206, 204}        % 優しいソフトローズ

% --- Headers & Footers (Strictly No text collision) ---
\pagestyle{fancy}
\fancyhf{}
\fancyhead[L]{\small\color{darkslate}\textbf{第\thechapter 講\quad \leftmark}}
\fancyhead[R]{\small\color{darkslate}デュードニー数学パズル入門}
\fancyfoot[C]{\color{darkslate}\thepage}
\renewcommand{\headrulewidth}{0.5pt}
\renewcommand{\chaptermark}[1]{\markboth{#1}{}}

% --- Hyperref ---
\usepackage{hyperref}
\hypersetup{
    colorlinks=true,
    linkcolor=primaryblue,
    citecolor=accentcrimson,
    urlcolor=primaryblue,
    pdftitle={古典で学び直すデュードニー数学パズル入門},
    pdfauthor={教育情報科学研究会}
}

\begin{document}

% --- Title Page ---
\begin{titlepage}
\centering
\vspace*{1.5cm}
{\color{primaryblue}\rule{\textwidth}{2pt}}\par
\vspace{1.0cm}
{\Huge\textbf{古典で学び直す}\par}
\vspace{0.4cm}
{\Huge\textbf{デュードニー数学パズル入門}\par}
\vspace{0.6cm}
{\Large\textbf{思考の幾何学・グラフ理論・数論の魔術}\par}
\vspace{0.8cm}
{\large —— 珠玉の430問から読み解く近代数学の核心 ——\par}
\vspace{1.0cm}
{\color{accentcrimson}\rule{0.6\textwidth}{1.5pt}}\par
\vspace{2.0cm}

\begin{tcolorbox}[colback=softblue,colframe=primaryblue,width=0.85\textwidth,arc=3mm,boxrule=1.2pt]
\centering\large
\textbf{ヘンリー・アーネスト・デュードニー} 原著\\
\textit{Amusements in Mathematics} (1917)\\
\vspace{0.5em}
\normalsize
50代、60代から始める、知的好奇心を満たす大人のための数学再入門
\end{tcolorbox}

\vfill
{\large 教育情報科学研究会 編著\par}
\vspace{1.0cm}
{\small 本書は英国の偉大なパズル作家H. E. デュードニーの名著『Amusements in Mathematics』を原典とし、\\現代数学（グラフ理論・線形代数・整数論・群論）の視点から厳選・再構成した解説書です。\par}
\vspace{1.0cm}
\end{titlepage}

% --- Front Matter ---
\frontmatter
\tableofcontents
\clearpage

% --- Introduction ---
\chapter*{発刊にあたって —— なぜいまデュードニーなのか}
\addcontentsline{toc}{chapter}{発刊にあたって —— なぜいまデュードニーなのか}

20世紀初頭のエドワード朝ロンドン。数学者たちが象牙の塔で格闘していた抽象理論を、市民の日常言語とユーモア溢れる物語に託して新聞・雑誌に連載し、英国中を熱狂させた一人の男がいました。**ヘンリー・アーネスト・デュードニー（Henry Ernest Dudeney, 1857--1930）**です。

彼が生み出した430問を超えるパズルを集大成した『Amusements in Mathematics』（1917年刊行）は、単なる暇つぶしのクイズ集ではありません。そこには、18世紀のオイラーによるグラフ理論の創始、ガウスの合同式による数論、さらには現代のコンピュータ・サイエンスを支えるアルゴリズムの基礎が、驚くほど直観的で美しい形で埋め込まれています。

大人の学び直しにおいて、数式を暗記する受験数学は不要です。必要なのは、**「なぜそのように考えるのか」「背後にある数学的な構造は何か」という本質への驚きと納得**です。

本書は、デュードニーの代表的傑作を8つの講義形式にまとめ、現代の視点からその数学的構造を徹底的に紐解きます。美しい図解とともに、知的な冒険の旅へ出かけましょう。

\clearpage

% --- Main Matter ---
\mainmatter

\chapter[ハバダッシャーの奇跡]{第1講：ハバダッシャーの奇跡 —— 正三角形から正方形への幾何学的変身}
\vspace{-1.0em}
\begin{center}
{\large\color{accentcrimson}\textbf{ヒンジ結合による図形分割・裁ち合わせ（Dissection）の極致}}
\end{center}
\vspace{1.0em}

幾何学において、ある多角形を有限個のピースに切り分け、並べ替えて別の多角形を隙間なく組み立て直す操作を「裁ち合わせ（Dissection / 等積変形）」と呼びます。

19世紀末、ハンガリーの数学者ボヤイ・ファルカシュとドイツのゲアウィンは、「面積の等しい任意の二つの多角形は、有限回の直線切断と並べ替えによって互いに移り変わることができる」というボヤイ・ゲアウィンの定理を証明しました。しかし、実際に「何個のピースで、どのように切断すればよいか」を見つけ出すことは、極めて高度な直観と数学的構成力を要求する難問です。

この分野における歴史上最大の傑作と称えられるのが、デュードニーが1902年に発表した「ハバダッシャー（洋品店主）のパズル（The Haberdasher's Puzzle, \#144・\#147の発展）」です。
彼は、一辺の長さが等しい正三角形を、わずか**「4つのピース」**に切り分け、それらがヒンジ（蝶番）で数珠つなぎになったまま回転させるだけで、完璧な「正方形」へと姿を変える驚異的な解法を発見しました。

正三角形の一辺を $a$ とすると面積は $S = \frac{\sqrt{3}}{4} a^2$ であり、これと等しい面積をもつ正方形の一辺は $x = \sqrt{S} = a \sqrt{\frac{\sqrt{3}}{4}} \approx 0.658 a$ となります。デュードニーは、正三角形の各辺の中点を基軸とし、幾何学的な垂線と中点連結を組み合わせることで、無理数 $\sqrt[4]{3}$ を見事に図形上で調和させ、4片の連鎖変形を実現したのです。

\vspace{1.2em}
\begin{tcolorbox}[colback=softblue,colframe=primaryblue,arc=2.5mm,boxrule=1.2pt,title={\textbf{重要ポイント：裁ち合わせの数理とヒンジ連鎖}}]
1. \textbf{ボヤイ・ゲアウィンの定理}：面積が等しい任意の2つの多角形は、有限個のピースへの分割によって必ず互いに変換可能である。
2. \textbf{最小ピース数の探求}：正三角形から正方形への変換は「4ピース」が理論的最小限界であることが証明されている。
3. \textbf{ヒンジ変形の幾何}：4つのピースは頂点のヒンジ（蝶番）でピン留めされており、裏返すことなく平面的にチェーン状に回すだけで三角形と正方形を往復する。
\end{tcolorbox}

\vspace{1.5em}
\begin{figure}[htbp]
\centering
\begin{tikzpicture}[>=Stealth, scale=1.05]
    % --- Left: Equilateral Triangle with 4 Dissected Pieces ---
    \begin{scope}[shift={(0,0)}]
        \coordinate (A) at (0, 3.464);
        \coordinate (B) at (-2, 0);
        \coordinate (C) at (2, 0);

        \coordinate (D) at (-1, 1.732);      % Midpoint AB (Hinge H1)
        \coordinate (E) at (1, 1.732);       % Midpoint AC (Hinge H2)
        \coordinate (G) at (1.0, 0);         % Midpoint BC (Hinge H3)
        \coordinate (F) at (0.0, 0);
        \coordinate (H) at (-0.5, 0.866);
        \coordinate (K) at (0.5, 0.866);

        % 4 Interlocking Pieces (Zero gap)
        \filldraw[fill=pieceA, draw=primaryblue, line width=1.3pt] (A) -- (E) -- (K) -- (D) -- cycle;
        \filldraw[fill=pieceB, draw=primaryblue, line width=1.3pt] (E) -- (C) -- (G) -- (K) -- cycle;
        \filldraw[fill=pieceC, draw=primaryblue, line width=1.3pt] (K) -- (G) -- (F) -- (H) -- cycle;
        \filldraw[fill=pieceD, draw=primaryblue, line width=1.3pt] (D) -- (H) -- (F) -- (B) -- cycle;

        \draw[line width=1.8pt, primaryblue] (A) -- (B) -- (C) -- cycle;

        % Hinges
        \filldraw[fill=warmgold, draw=darkslate, line width=1.2pt] (D) circle (3.5pt) node[left=4pt, font=\small\bfseries] {$H_1$};
        \filldraw[fill=warmgold, draw=darkslate, line width=1.2pt] (E) circle (3.5pt) node[right=4pt, font=\small\bfseries] {$H_2$};
        \filldraw[fill=warmgold, draw=darkslate, line width=1.2pt] (G) circle (3.5pt) node[below=4pt, font=\small\bfseries] {$H_3$};

        \node[above=2pt] at (A) {\textbf{\large $A$}};
        \node[below left=2pt] at (B) {\textbf{\large $B$}};
        \node[below right=2pt] at (C) {\textbf{\large $C$}};

        \node[below=16pt, align=center] at (0, 0) {
            \textbf{【正三角形】}\quad 面積 $S = \frac{\sqrt{3}}{4} a^2$\\[3pt]
            \footnotesize 3箇所のヒンジ（$H_1, H_2, H_3$）で鎖状連結
        };
    \end{scope}

    % --- Center: Transformation Arrow ---
    \draw[->, line width=2.4pt, accentcrimson] (2.4, 1.732) -- (4.3, 1.732)
        node[midway, above=4pt, align=center] {\textbf{\small ヒンジ連鎖回転}}
        node[midway, below=4pt, align=center] {\scriptsize\color{darkslate} 平面的に開閉変形};

    % --- Right: Reassembled Perfect Square ---
    \begin{scope}[shift={(5.9, 0.42)}]
        \def\sq{2.632}

        \coordinate (S1) at (0, 0);
        \coordinate (S2) at (\sq, 0);
        \coordinate (S3) at (\sq, \sq);
        \coordinate (S4) at (0, \sq);

        \coordinate (SP1) at (1.0, 0);
        \coordinate (SP2) at (\sq, 1.0);
        \coordinate (SP3) at (\sq - 1.0, \sq);
        \coordinate (SP4) at (0, \sq - 1.0);
        \coordinate (SC)  at (1.316, 1.316);

        \filldraw[fill=pieceA, draw=primaryblue, line width=1.3pt] (S4) -- (SP3) -- (SC) -- (SP4) -- cycle;
        \filldraw[fill=pieceB, draw=primaryblue, line width=1.3pt] (SP3) -- (S3) -- (SP2) -- (SC) -- cycle;
        \filldraw[fill=pieceC, draw=primaryblue, line width=1.3pt] (SC) -- (SP2) -- (S2) -- (SP1) -- cycle;
        \filldraw[fill=pieceD, draw=primaryblue, line width=1.3pt] (SP4) -- (SC) -- (SP1) -- (S1) -- cycle;

        \draw[line width=1.8pt, accentcrimson] (S1) -- (S2) -- (S3) -- (S4) -- cycle;

        \filldraw[fill=warmgold, draw=darkslate, line width=1.2pt] (SP4) circle (3.5pt) node[left=4pt, font=\small\bfseries] {$H_1$};
        \filldraw[fill=warmgold, draw=darkslate, line width=1.2pt] (SP3) circle (3.5pt) node[above=4pt, font=\small\bfseries] {$H_2$};
        \filldraw[fill=warmgold, draw=darkslate, line width=1.2pt] (SP1) circle (3.5pt) node[below=4pt, font=\small\bfseries] {$H_3$};

        \node[below=16pt, align=center] at (\sq/2, 0) {
            \textbf{【完全な正方形】}\quad 一辺 $x = a \sqrt{\frac{\sqrt{3}}{4}}$\\[3pt]
            \footnotesize 同色ピースが隙間なく正方形を再構成
        };
    \end{scope}
\end{tikzpicture}
\caption{ハバダッシャー・パズル：正三角形から正方形への4分割ヒンジ連鎖変形。}
\end{figure}

\clearpage
\chapter[オイラーの一筆書き]{第2講：オイラーの一筆書き定理 —— ケーニヒスベルクから始まるグラフ理論}
\vspace{-1.0em}
\begin{center}
{\large\color{accentcrimson}\textbf{頂点の次数（偶奇性）が握るネットワーク走破の運命}}
\end{center}
\vspace{1.0em}

1736年、大数学者レオンハルト・オイラーは「ケーニヒスベルクの7つの橋」問題を解決し、現代のトポロジー（位相幾何学）およびグラフ理論の扉を開きました。この理論を誰でも親しめる魅力的なパズルへと昇華させたのがデュードニーです（\#220--\#261）。

一筆書きができるかどうかは、図形の大きさや線の長さとは一切関係がありません。問題となるのは、**「各頂点（交差点）に集まる線の本数（次数 / degree）」**だけです。

線の本数が奇数である頂点を「奇頂点（Odd Vertex）」、偶数である頂点を「偶頂点（Even Vertex）」と呼びます。1本の線を描いて通過するたびに、その交差点では「入る線」と「出る線」のペア（2本）が消費されます。

したがって：
1. **奇頂点が0個**の場合：どの点から出発しても一筆書きが可能で、必ず出発点に戻る（オイラー閉路）。
2. **奇頂点が2個**の場合：一方の奇頂点から出発し、もう一方の奇頂点で終わる一筆書きが可能（オイラー路）。
3. **奇頂点が4個以上**の場合：いかなる工夫を凝らしても一筆書きは**数学的に不可能**である。

デュードニーはこのオイラーの定理を駆使し、鉄道の引き込み線パズルや家々の巡回ルート問題など、数々の魅力的な傑作を生み出しました。

\vspace{1.2em}
\begin{tcolorbox}[colback=softblue,colframe=primaryblue,arc=2.5mm,boxrule=1.2pt,title={\textbf{重要ポイント：オイラーの一筆書き定理}}]
1. \textbf{幾何学的寸法の捨象}：線の曲がりや長さは無関係であり、頂点の接続関係（次数）のみが本質を決める。
2. \textbf{奇頂点数の偶数性}：1本の線は2つの端点を持つため、グラフ全体の奇頂点の総数は常に偶数（0, 2, 4, 6...）個である。
3. \textbf{一筆書きの完全判定}：奇頂点が0個または2個の連結グラフのみが一筆書き可能であり、4個以上は絶対に描けない。
\end{tcolorbox}

\vspace{1.5em}
\begin{figure}[htbp]
\centering
\begin{tikzpicture}[>=Stealth, scale=1.1]
    % Case A: Euler Tour (All Even Vertices)
    \begin{scope}[shift={(0,0)}]
        \node[draw=primaryblue, line width=1.5pt, fill=softblue, circle, inner sep=4pt] (A1) at (0, 2.2) {\textbf{4}};
        \node[draw=primaryblue, line width=1.5pt, fill=softblue, circle, inner sep=4pt] (B1) at (-1.8, 0) {\textbf{2}};
        \node[draw=primaryblue, line width=1.5pt, fill=softblue, circle, inner sep=4pt] (C1) at (1.8, 0) {\textbf{2}};
        \node[draw=primaryblue, line width=1.5pt, fill=softblue, circle, inner sep=4pt] (D1) at (0, -1.2) {\textbf{4}};

        \draw[line width=1.4pt, primaryblue] (A1) -- (B1) -- (D1) -- (C1) -- (A1);
        \draw[line width=1.4pt, primaryblue] (A1) -- (D1);
        \draw[line width=1.4pt, primaryblue] (-1.8,0) to[bend left=30] (1.8,0);

        \node[below=14pt, align=center] at (0, -1.2) {
            \textbf{【奇頂点 0個】}\quad 全て偶数次\\
            \small どの頂点からでも出発点に戻る一筆書きが可能
        };
    \end{scope}

    % Case B: Non-Eulerian (4 Odd Vertices)
    \begin{scope}[shift={(7.5, 0)}]
        \node[draw=accentcrimson, line width=1.5pt, fill=accentcrimson!15, circle, inner sep=4pt] (A2) at (-1.5, 2.0) {\textbf{3}};
        \node[draw=accentcrimson, line width=1.5pt, fill=accentcrimson!15, circle, inner sep=4pt] (B2) at (1.5, 2.0) {\textbf{3}};
        \node[draw=accentcrimson, line width=1.5pt, fill=accentcrimson!15, circle, inner sep=4pt] (C2) at (-1.5, -0.6) {\textbf{3}};
        \node[draw=accentcrimson, line width=1.5pt, fill=accentcrimson!15, circle, inner sep=4pt] (D2) at (1.5, -0.6) {\textbf{3}};

        \draw[line width=1.4pt, darkslate] (A2) -- (B2) -- (D2) -- (C2) -- (A2);
        \draw[line width=1.4pt, darkslate] (A2) -- (D2);
        \draw[line width=1.4pt, darkslate] (B2) -- (C2);

        \node[below=14pt, align=center] at (0, -0.6) {
            \textbf{【奇頂点 4個】}\quad 完全グラフ $K_4$\\
            \small 奇頂点が4個 $\to$ \textbf{一筆書きは数学的に不可能}
        };
    \end{scope}
\end{tikzpicture}
\caption{オイラーの一筆書き判定原理。各頂点に集まる線の本数（偶数・奇数）が運命を完全に決定する。}
\end{figure}

\clearpage
\chapter[チェス盤の数学]{第3講：チェス盤の数学 —— ナイト・ツアーと8クイーン問題}
\vspace{-1.0em}
\begin{center}
{\large\color{accentcrimson}\textbf{二次元格子上の探索・対称群・ハミルトン路の数理}}
\end{center}
\vspace{1.0em}

64個の白黒マスからなるチェス盤は、それ自体が豊かな離散数学の実験場です（\#288--\#361）。デュードニーはチェス盤を用いた数学パズルの第一人者でもありました。

代表的な古典問題の一つが「ナイト・ツアー（Knight's Tour, \#339--\#340）」です。チェスのナイト（八方に「2マス進んで1マス曲がる」L字型に跳ぶ駒）を盤上に置き、64マスすべてのマスを重複することなくちょうど1度ずつ訪問して元の位置に戻る（または終わる）ことができるか、という問題です。

グラフ理論の言葉で言えば、これは「ナイトの移動可能なマス同士を結んだグラフにおいて**ハミルトン閉路（Hamiltonian Cycle）**が存在するか」という問題に他なりません。

ここで極めて重要な数学的性質が**「色の交代性（パリティ）」**です。ナイトはどのような跳び方をしても、**白マスから黒マスへ、黒マスから白マスへと、1手ごとに必ずマスの色が反転**します。したがって、ナイト・ツアーが完結するためには、白32マス・黒32マスを厳密に交互に訪れなければならないという絶対的な制約（二部グラフの性質）が課されるのです。

\vspace{1.2em}
\begin{tcolorbox}[colback=softblue,colframe=primaryblue,arc=2.5mm,boxrule=1.2pt,title={\textbf{重要ポイント：二次元格子探索とパリティ不変量}}]
1. \textbf{二部グラフ構造}：チェス盤の白マスと黒マスは完全な二部グラフをなし、ナイトの移動は異なるグループ間の辺のみを通る。
2. \textbf{ハミルトン路の存在条件}：全64マスを踏破するには、奇数手目は白、偶数手目は黒（またはその逆）という厳密な交代性が必須となる。
3. \textbf{局所対称性と大域探索}：コーナー（角）のマスは進入路が2本しか存在しないため、解法の決定において最優先の制約となる。
\end{tcolorbox}

\vspace{1.5em}
\begin{figure}[htbp]
\centering
\begin{tikzpicture}[scale=0.9, >=Stealth]
    % 5x5 Chessboard with checkered pattern
    \begin{scope}[shift={(0,0)}]
        \foreach \x in {-2,-1,0,1,2} {
            \foreach \y in {-2,-1,0,1,2} {
                \pgfmathtruncatemacro{\parity}{mod(\x + \y + 10, 2)}
                \ifnum\parity=0
                    \fill[white] (\x-0.5, \y-0.5) rectangle (\x+0.5, \y+0.5);
                \else
                    \fill[softblue!80!primaryblue!22] (\x-0.5, \y-0.5) rectangle (\x+0.5, \y+0.5);
                \fi
                \draw[darkslate!35, thin] (\x-0.5, \y-0.5) rectangle (\x+0.5, \y+0.5);
            }
        }
        \draw[line width=1.6pt, darkslate] (-2.5, -2.5) rectangle (2.5, 2.5);

        % Center Knight at (0,0) [White Square]
        \filldraw[fill=primaryblue, draw=primaryblue] (0, 0) circle (0.38);
        \node[white] at (0, 0) {\textbf{\large N}};

        % 8 Target Squares (All colored squares by chess geometry!)
        \foreach \tx/\ty in {1/2, -1/2, 1/-2, -1/-2, 2/1, -2/1, 2/-1, -2/-1} {
            \fill[accentcrimson!25] (\tx-0.5, \ty-0.5) rectangle (\tx+0.5, \ty+0.5);
            \draw[accentcrimson, line width=1.2pt] (\tx-0.5, \ty-0.5) rectangle (\tx+0.5, \ty+0.5);
            \filldraw[fill=accentcrimson, draw=white, line width=0.8pt] (\tx, \ty) circle (3pt);

            % L-shaped trajectory lines
            \ifnum\tx=1 \ifnum\ty=2 \draw[->, line width=1.5pt, accentcrimson] (0,0) -- (0, 2) -- (1, 2); \fi\fi
            \ifnum\tx=-1 \ifnum\ty=2 \draw[->, line width=1.1pt, accentcrimson, dashed] (0,0) -- (0, 2) -- (-1, 2); \fi\fi
            \ifnum\tx=2 \ifnum\ty=1 \draw[->, line width=1.5pt, accentcrimson] (0,0) -- (2, 0) -- (2, 1); \fi\fi
            \ifnum\tx=-2 \ifnum\ty=1 \draw[->, line width=1.1pt, accentcrimson, dashed] (0,0) -- (-2, 0) -- (-2, 1); \fi\fi
            \ifnum\tx=1 \ifnum\ty=-2 \draw[->, line width=1.1pt, accentcrimson, dashed] (0,0) -- (0, -2) -- (1, -2); \fi\fi
            \ifnum\tx=-1 \ifnum\ty=-2 \draw[->, line width=1.1pt, accentcrimson, dashed] (0,0) -- (0, -2) -- (-1, -2); \fi\fi
            \ifnum\tx=2 \ifnum\ty=-1 \draw[->, line width=1.1pt, accentcrimson, dashed] (0,0) -- (2, 0) -- (2, -1); \fi\fi
            \ifnum\tx=-2 \ifnum\ty=-1 \draw[->, line width=1.1pt, accentcrimson, dashed] (0,0) -- (-2, 0) -- (-2, -1); \fi\fi
        }

        \node[below=14pt, align=center] at (0, -2.5) {
            \small\textbf{ナイトの現在地（白マス中央）}\\
            \footnotesize 八方の赤マス（すべて黒マス）へ到達可能
        };
    \end{scope}

    % Right Explanation Box (Adequate spacing)
    \begin{scope}[shift={(6.7, 0)}]
        \node[draw=primaryblue, line width=1.2pt, fill=softblue, rounded corners=3mm, inner sep=9pt, text width=6.0cm, align=left] at (0, 0) {
            \textbf{\large 【ナイト移動の不変量】}\\[6pt]
            \small
            \textbf{・移動ベクトル}：$(\pm 1, \pm 2)$ または $(\pm 2, \pm 1)$\\[3pt]
            \textbf{・格子点間距離}：$\sqrt{1^2 + 2^2} = \sqrt{5}$\\[3pt]
            \textbf{・マスの色交代性}：\\
            \quad \colorbox{white}{\fbox{\strut 白マス}} $\longrightarrow$ \colorbox{accentcrimson!20}{\fbox{\strut\bfseries 常に黒マス}}\\[4pt]
            1手進むごとにマスの色が必ず反転する！\\[6pt]
            \textbf{・64手完全巡回} $\implies$ ハミルトン閉路\\
            対称性とパリティ（偶奇性）が全マス踏破の成否を握る。
        };
    \end{scope}
\end{tikzpicture}
\caption{チェスのナイトがもつ「L字移動」の対称性と白黒反転パリティ。ナイト・ツアーの数学的基礎。}
\end{figure}

\clearpage
\chapter[川渡りのアルゴリズム]{第4講：川渡りパズルの状態空間探索 —— 狼と山羊とキャベツの論理}
\vspace{-1.0em}
\begin{center}
{\large\color{accentcrimson}\textbf{制約充足問題（CSP）と有向グラフ最短経路アルゴリズム}}
\end{center}
\vspace{1.0em}

「農夫が狼、山羊、キャベツを一艘の小舟で対岸へ運ばなければならない。舟には農夫のほかに1つしか乗せられない。農夫がいないと、狼は山羊を食べ、山羊はキャベツを食べてしまう。どう運べば全員無事に対岸へ渡れるか？」（\#373--\#377）。

中世から伝わるこの有名な古典パズルは、現代情報科学における**「状態空間探索（State Space Search）」**および**「制約充足問題（Constraint Satisfaction Problem）」**の原点です。

各時刻において、「誰が左岸にいて、誰が右岸にいるか」という状態は全部で $2^4 = 16$ 通り存在します。しかし、「農夫不在で狼と山羊が同居」「農夫不在で山羊とキャベツが同居」という禁止状態（制約違反）を除外すると、許される安全な状態はわずか**10通り**に絞られます。

この安全な状態をノード（節点）、舟による合法的な移動をエッジ（枝）とする「有向グラフ」を描けば、パズルを解くことは**「初期状態から目標状態への最短経路を探索する問題」**へと完全に帰着されます。この問題のクライマックスは、4手目に「一旦対岸へ運んだ山羊を連れ戻す」という逆転の発想です。局所的な前進にとらわれず、状態空間全体を俯瞰する視点こそが数学的思考の本質です。

\vspace{1.2em}
\begin{tcolorbox}[colback=softblue,colframe=primaryblue,arc=2.5mm,boxrule=1.2pt,title={\textbf{重要ポイント：状態空間と幅優先探索（BFS）}}]
1. \textbf{状態のモデル化}：現実の制約をBoolean（真偽値）やビット列として定式化する。
2. \textbf{制約による枝刈り}：不可能な状態を探索木から即座に排除し、探索空間を劇的に縮小させる。
3. \textbf{一時的後退の価値}：最短解に至る経路には、一見無駄に見える「元の岸へ戻す」ステップが不可欠となる。
\end{tcolorbox}

\vspace{1.5em}
\begin{figure}[htbp]
\centering
\begin{tikzpicture}[>=Stealth, scale=0.92]
    % State nodes (Arranged neatly with compact multi-line text)
    \begin{scope}[shift={(0,0)}]
        \node[draw=primaryblue, line width=1.4pt, fill=softblue, rounded corners=2mm, inner sep=6pt] (S0) at (0, 3.2) {\textbf{初期：左岸(人, 狼, 羊, 甘)}};
        \node[draw=darkslate, line width=1.1pt, fill=white, rounded corners=2mm, inner sep=5pt] (S1) at (0, 1.8) {\small 羊を対岸へ渡す};
        \node[draw=darkslate, line width=1.1pt, fill=white, rounded corners=2mm, inner sep=5pt] (S2) at (0, 0.4) {\small 人が単身で戻る};
        \node[draw=darkslate, line width=1.1pt, fill=white, rounded corners=2mm, inner sep=5pt, align=center] (S3) at (-2.3, -1.0) {\small 狼を渡し\\[2pt]\small 羊を連れ戻す};
        \node[draw=darkslate, line width=1.1pt, fill=white, rounded corners=2mm, inner sep=5pt, align=center] (S4) at (2.3, -1.0) {\small 甘を渡し\\[2pt]\small 羊を連れ戻す};
        \node[draw=darkslate, line width=1.1pt, fill=white, rounded corners=2mm, inner sep=5pt] (S5) at (0, -2.4) {\small 甘（または狼）を渡す};
        \node[draw=accentcrimson, line width=1.6pt, fill=accentcrimson!15, rounded corners=2mm, inner sep=6pt] (S6) at (0, -3.8) {\textbf{達成：全員無事に対岸へ！}};

        \draw[->, line width=1.3pt, primaryblue] (S0) -- (S1);
        \draw[->, line width=1.3pt, primaryblue] (S1) -- (S2);
        \draw[->, line width=1.3pt, primaryblue] (S2) -- (S3);
        \draw[->, line width=1.3pt, primaryblue] (S2) -- (S4);
        \draw[->, line width=1.3pt, primaryblue] (S3) -- (S5);
        \draw[->, line width=1.3pt, primaryblue] (S4) -- (S5);
        \draw[->, line width=1.6pt, accentcrimson] (S5) -- (S6);
    \end{scope}

    % Right Key insight box (Shifted further right to guarantee zero collision)
    \begin{scope}[shift={(7.3, -0.2)}]
        \node[draw=warmgold!80!black, fill=warmgold!10, rounded corners=2mm, text width=4.6cm, inner sep=8pt, align=left] at (0, 0) {
            \textbf{\small 【逆転の定石】}\\[4pt]
            \footnotesize
            荷物を対岸に運ぶだけでなく、一時的に「羊を連れ帰る」という一見無駄な後退が、袋小路の唯一の突破口となる。\\
            \vspace{4pt}
            状態空間グラフにおける最短路アルゴリズムの真骨頂。
        };
    \end{scope}
\end{tikzpicture}
\caption{川渡りパズルの状態遷移ネットワーク。禁忌状態を排除した合法状態空間の探索グラフ。}
\end{figure}

\clearpage
\chapter[魔方陣の数理]{第5講：魔方陣の代数構造 —— 線形代数と対称群が織りなす宇宙}
\vspace{-1.0em}
\begin{center}
{\large\color{accentcrimson}\textbf{古代中国の「洛書」からアフィン空間への昇華}}
\end{center}
\vspace{1.0em}

縦・横・対角線の数の和がすべて等しくなるように正方形に数を並べた「魔方陣（Magic Square, \#399--\#412）」。古代中国の易経に登場する「洛書（3次魔方陣）」をはじめ、西洋ではデューラーの版画『メランコリア I』に描かれるなど、人類を魅了し続けてきた神秘の数配置です。

デュードニーはこの魔方陣を、単なる試行錯誤のパズルから厳密な代数学の対象へと引き上げました。

最も基本的な $3 \times 3$ の魔方陣（1から9までの数字）において、全数字の総和は $1 + 2 + \dots + 9 = 45$ です。3行あるため、各行・各列の和（魔方陣定数）は必ず $45 \div 3 = 15$ にならなければなりません。

さらに、中央のマスを通るラインは4本（行1・列1・対角線2）あり、これらを連立させることで、**「中央の数は必ず 5 でなければならない」**ことが数学的に一意に証明されます。

さらに進んで、現代の線形代数では、魔方陣全体が「連立一次方程式の解空間（アフィン部分空間）」をなすことが知られています。自由度（次元）を計算することで、基本となる独立な魔方陣の線形結合としてすべての魔方陣を生成できるのです。

\vspace{1.2em}
\begin{tcolorbox}[colback=softblue,colframe=primaryblue,arc=2.5mm,boxrule=1.2pt,title={\textbf{重要ポイント：魔方陣の不変量と線形構造}}]
1. \textbf{定数の必然性}：$n$ 次魔方陣の行和・列和は、総和の対称性から $M = \frac{n(n^2+1)}{2}$ と厳密に一意決定される。
2. \textbf{中心数の定理}：奇数次魔方陣の中心マスには、必ず全数字の中央値（3次ならば5）が配置される。
3. \textbf{対称変換群}：回転（4通り）と反転（2通り）による8個の変換は、二面体群 $D_4$ を構成し、本質的な解は1種類に帰着される。
\end{tcolorbox}

\vspace{1.5em}
\begin{figure}[htbp]
\centering
\begin{tikzpicture}[>=Stealth, scale=1.1]
    % 3x3 Magic Square
    \begin{scope}[shift={(0,0)}]
        \draw[step=1.2cm, darkslate, line width=1.4pt] (0,0) grid (3.6,3.6);
        \draw[line width=2.0pt, primaryblue] (0,0) rectangle (3.6,3.6);

        % Numbers
        \node at (0.6, 3.0) {\textbf{\LARGE 8}};
        \node at (1.8, 3.0) {\textbf{\LARGE 1}};
        \node at (3.0, 3.0) {\textbf{\LARGE 6}};

        \node at (0.6, 1.8) {\textbf{\LARGE 3}};
        \node[accentcrimson] at (1.8, 1.8) {\textbf{\LARGE 5}};
        \node at (3.0, 1.8) {\textbf{\LARGE 7}};

        \node at (0.6, 0.6) {\textbf{\LARGE 4}};
        \node at (1.8, 0.6) {\textbf{\LARGE 9}};
        \node at (3.0, 0.6) {\textbf{\LARGE 2}};

        \node[below=14pt, align=center] at (1.8, 0) {
            \textbf{【3次魔方陣（洛書）】}\\
            \small 各行・列・対角線の和 $= \mathbf{15}$（中心は常に $5$）
        };
    \end{scope}

    % Mathematical Decomposition Box
    \node[draw=primaryblue, fill=softblue, rounded corners=2mm, text width=6.2cm, inner sep=9pt, align=left] at (7.2, 1.8) {
        \textbf{\large 【魔方陣の数学的構造】}\\[5pt]
        \small 任意の3次魔方陣は、中心数 $x$、パラメータ $a, b$ を用いる線形モデルとして表現：\\[4pt]
        \centering
        \begin{tabular}{|c|c|c|}
        \hline
        $x+b$ & $x-a-b$ & $x+a$ \\ \hline
        $x+a-b$ & $x$ & $x-a+b$ \\ \hline
        $x-a$ & $x+a+b$ & $x-b$ \\ \hline
        \end{tabular}\\[6pt]
        \raggedright\footnotesize
        線形代数における3次元ベクトル空間の美しい具現化。
    };
\end{tikzpicture}
\caption{3次魔方陣の代数的構造。中心数「5」の決定性と線形空間としての一般解。}
\end{figure}

\clearpage
\chapter[迷路の幾何学]{第6講：迷路の幾何学と脱出アルゴリズム —— トレモーの解法}
\vspace{-1.0em}
\begin{center}
{\large\color{accentcrimson}\textbf{トポロジーと深さ優先探索（DFS）の先駆的アルゴリズム}}
\end{center}
\vspace{1.0em}

入り組んだ通路の迷宮から生還するための数学的手法とは何か？ デュードニーの第13章「迷路の問題（Mazes and How to Thread Them）」は、古今東西の有名な迷路（ハンプトンコート宮殿の生垣迷路など）を取り上げ、その幾何学的本質を解き明かしています。

初等的な迷路では「壁に右手を当てて歩き続ければ必ず脱出できる（右手法 / 壁づたい法）」と言われます。しかし、これは「迷路が単連結（環状の孤島がない）」場合にしか通用しません。中央に独立したループ構造（多重連結迷路）がある場合、右手法では同じ環状路を無限に周回してしまう罠に陥ります。

いかなる複雑な迷路であっても確実に脱出できる完全なアルゴリズムとして知られるのが、19世紀のフランスの数学者シャルル・ピエール・トレモーが考案した**「トレモーのアルゴリズム（Trémaux's Algorithm）」**です。

ルールは極めてシンプルです：
1. 通路を初めて進むときは、床に「1本線」の印をつける。
2. もし行き止まりに突き当たったら、回れ右をして戻り、通路に「2本目」の印をつけて**その通路を封鎖（二重印）**する。
3. すでに印がついている交差点に別の道から出た場合、決してその道には進まず、来た道を引き返して2本目の印をつける。

この操作は、コンピュータ・アルゴリズムにおける**「深さ優先探索（Depth-First Search, DFS）」そのもの**です。有限の迷路において、各通路は最大でも「往路」と「復路」の計2回しか通過されず、無限ループに陥ることなく確実に脱出路を発見できることが数学的に証明されています。

\vspace{1.2em}
\begin{tcolorbox}[colback=softblue,colframe=primaryblue,arc=2.5mm,boxrule=1.2pt,title={\textbf{重要ポイント：トレモーのアルゴリズムとDFS}}]
1. \textbf{右手法の限界}：島（ループ）を含む多重連結迷路では、壁づたい法は無限周回トラップに囚われる。
2. \textbf{印による状態記憶}：未踏（0本）、探索中（1本）、走破・袋小路（2本）という3状態の記録が完全走破を保証する。
3. \textbf{計算量の最適性}：通路の総数を $E$ とすると、最大移動回数は高々 $2E$ 回で終了し、線形時間 $O(E)$ でゴールに到達する。
\end{tcolorbox}

\vspace{1.5em}
\begin{figure}[htbp]
\centering
\begin{tikzpicture}[>=Stealth, scale=0.95]
    % Maze graph nodes (Wide, spacious layout with guaranteed zero overlap)
    \node[draw=primaryblue, line width=1.6pt, fill=softblue, circle, inner sep=5pt] (IN) at (0, 0) {\textbf{入口}};
    \node[draw=darkslate, line width=1.3pt, fill=white, circle, inner sep=5pt] (A) at (3.8, 1.6) {\textbf{分岐 $A$}};
    \node[draw=accentcrimson, line width=1.5pt, fill=accentcrimson!15, rectangle, rounded corners=2mm, inner sep=6pt] (DEAD) at (9.4, 1.6) {\textbf{行き止まり}};
    \node[draw=darkslate, line width=1.3pt, fill=white, circle, inner sep=5pt] (B) at (3.8, -1.4) {\textbf{分岐 $B$}};
    \node[draw=forestgreen, line width=1.6pt, fill=forestgreen!20, circle, inner sep=6pt] (OUT) at (9.4, -1.4) {\textbf{出口}};

    % Path ①: IN -> A
    \draw[->, line width=1.6pt, primaryblue] (IN) -- (A)
        node[midway, above left=2pt] {\small ① 進入（1本線）};

    % Path ②: A -> DEAD (Forward straight exploration)
    \draw[->, line width=1.6pt, darkslate] (A) -- (DEAD)
        node[midway, below=3pt] {\small ② 直進探索（1本線）};

    % Path ③: DEAD -> A (Upward curving backtrack: 100% separated from Path ④!)
    \draw[->, line width=1.6pt, accentcrimson, dashed] (DEAD) to[bend right=42]
        node[midway, above=4pt, accentcrimson] {\small ③ 二重印で封鎖して引き返す} (A);

    % Path ④: A -> B (Downward to unexplored path)
    \draw[->, line width=1.6pt, primaryblue] (A) -- (B)
        node[midway, right=6pt] {\small ④ 未踏路へ進行};

    % Path ⑤: B -> OUT (Goal)
    \draw[->, line width=2.0pt, forestgreen] (B) -- (OUT)
        node[midway, above=4pt, forestgreen] {\textbf{\small ⑤ ゴール到達！}};

    % Bottom Explanation Box with ample top margin
    \node[draw=primaryblue, line width=1.2pt, fill=white, rounded corners=2mm, text width=12.8cm, inner sep=8pt, align=left] at (4.7, -3.2) {
        \small \textbf{トレモーの規則}：通路に印をつけることで、「未踏の道」と「探索済みの袋小路」を厳密に区別する。\\
        いかなる多重連結迷路も、各通路を往復（計2回）以上通過することなく、有限時間で完全走破される。
    };
\end{tikzpicture}
\caption{トレモーの迷路走破アルゴリズム。袋小路の二重印封鎖と深さ優先探索のメカニズム。}
\end{figure}

\clearpage
\chapter[不定方程式と整数論]{第7講：ディオファントス方程式と整数論 —— 金銭・年齢・分配パズル}
\vspace{-1.0em}
\begin{center}
{\large\color{accentcrimson}\textbf{「解が整数でなければならない」という最も美しい制約の解法}}
\end{center}
\vspace{1.0em}

代数学の通常の方程式では、未知数の数よりも方程式の数が少ない場合、解は無数に存在して定まりません。しかし、「未知数はすべて正の整数（自然数）でなければならない」という条件がついた瞬間、状況は一変します。

このような方程式を、古代アレクサンドリアの数学者にちなんで**「不定方程式（ディオファントス方程式 / Diophantine Equation）」**と呼びます。

デュードニーの第1章「算術と代数の問題（\#1--\#141）」の多くは、このディオファントス方程式の鮮やかな応用です。
例えば、問題第1問「郵便局の困惑（\#1）」では、5シリング（60ペンス）の硬貨を出し、「2ペンス切手を何枚か、その6倍の枚数の1ペンス切手をくれ、そして残りを2ペンス半（2.5ペンス）切手でぴったり支払ってほしい」と求めます。

2ペンス切手の枚数を $x$ とすると、1ペンス切手は $6x$ 枚となります。2.5ペンス切手の枚数を $y$ と置くと、合計金額は：
\[
2x + 1(6x) + 2.5y = 60 \implies 8x + \frac{5}{2}y = 60 \implies 16x + 5y = 120
\]
この方程式において、$x, y$ はともに正の整数でなければなりません。
合同式（mod 5）をとると：
\[
16x \equiv 120 \pmod 5 \implies x \equiv 0 \pmod 5
\]
すなわち、$x$ は 5 の倍数でなければなりません。
$x = 5$ を代入すると、$16(5) + 5y = 80 + 5y = 120 \implies 5y = 40 \implies y = 8$。
もし $x \ge 10$ とすると $16x \ge 160 > 120$ となり、$y$ が負になってしまうため不適です。

したがって、可能な解は**「$x = 5, y = 8$」のただ1組**に確定します！
（2ペンス切手5枚、1ペンス切手30枚、2.5ペンス切手8枚）。

学校の数学では「連立方程式は式の数が足りなければ解けない」と教わりますが、現実の世界（切手の枚数や人数）がもつ「整数」という制約が、見事に唯一無二の答えを浮き彫りにするのです。

\vspace{1.2em}
\begin{tcolorbox}[colback=softblue,colframe=primaryblue,arc=2.5mm,boxrule=1.2pt,title={\textbf{重要ポイント：一次不定方程式と整数の剰余類}}]
1. \textbf{ベズーの等式}：方程式 $ax + by = c$ が整数解をもつための必要十分条件は、$c$ が $\gcd(a, b)$ の倍数であること。
2. \textbf{合同式（モジュロ演算）の威力}：法（modulus）をとることで変数を1つ消去し、解の候補を瞬時に絞り込む。
3. \textbf{不等式による境界画定}：非負整数条件（$x \ge 1, y \ge 1$）によって、無限に広がる解の中から現実の有限解を特定する。
\end{tcolorbox}

\vspace{1.5em}
\begin{figure}[htbp]
\centering
\begin{tikzpicture}[scale=0.95, >=Stealth]
    % Accurate coordinate system & grid
    \def\xscale{1.15}
    \def\yscale{0.22}

    % Grid lines (integer lattice)
    \foreach \x in {1,2,3,4,5,6,7,8} {
        \draw[gray!20, thin] (\x*\xscale, 0) -- (\x*\xscale, 26*\yscale);
        \node[below=3pt, font=\footnotesize\color{darkslate}] at (\x*\xscale, 0) {\x};
    }
    \foreach \y in {5,8,10,15,20,24} {
        \draw[gray!20, thin] (0, \y*\yscale) -- (8*\xscale, \y*\yscale);
        \node[left=3pt, font=\footnotesize\color{darkslate}] at (0, \y*\yscale) {\y};
    }

    % Origin label (single, clean)
    \node[below left=3pt, font=\footnotesize\color{darkslate}] at (0, 0) {0};

    % Axes with comfortable margins
    \draw[->, line width=1.3pt, darkslate] (-0.3, 0) -- (9.8*\xscale, 0)
        node[below right=2pt, font=\small\color{darkslate}] {\textbf{$x$（2ペンス切手基準枚数）}};
    \draw[->, line width=1.3pt, darkslate] (0, -0.3) -- (0, 28*\yscale)
        node[above=6pt, font=\small\color{darkslate}] {\textbf{$y$（2.5ペンス切手枚数）}};

    % Line: 16x + 5y = 120 (y = 24 - 3.2x)
    \draw[line width=2.0pt, primaryblue] (0, 24*\yscale) -- (7.5*\xscale, 0);
    \node[above right=3pt, primaryblue] at (1.8*\xscale, 18.2*\yscale) {\textbf{直線 $16x + 5y = 120$}};

    % Boundary points
    \filldraw[fill=white, draw=darkslate, line width=1.5pt] (0, 24*\yscale) circle (3.5pt)
        node[above right=3pt, font=\footnotesize, fill=white, inner sep=1.5pt] {$(0, 24)$ [切手無指定・不適]};

    \filldraw[fill=white, draw=darkslate, line width=1.5pt] (7.5*\xscale, 0) circle (3.5pt)
        node[above right=3pt, font=\footnotesize, fill=white, inner sep=1.5pt] {$(7.5, 0)$ [非整数]};

    % Non-integer points on the line for integer x
    \foreach \x/\y in {1/20.8, 2/17.6, 3/14.4, 4/11.2, 6/4.8, 7/1.6} {
        \filldraw[fill=gray!60, draw=darkslate] (\x*\xscale, \y*\yscale) circle (2.5pt);
        \draw[gray!45, dotted] (\x*\xscale, 0) -- (\x*\xscale, \y*\yscale);
    }

    % The Unique Integer Solution (5, 8)
    \draw[accentcrimson, line width=1.3pt, dashed] (5*\xscale, 0) -- (5*\xscale, 8*\yscale) -- (0, 8*\yscale);
    \filldraw[fill=accentcrimson, draw=primaryblue, line width=1.8pt] (5*\xscale, 8*\yscale) circle (5pt);
    \node[above right=8pt, fill=softblue, draw=accentcrimson, line width=1.2pt, rounded corners=2mm, inner sep=5pt] at (5*\xscale, 8*\yscale) {
        \textbf{\large 唯一の自然数解 $(x=5, y=8)$}
    };
\end{tikzpicture}
\caption{一次不定方程式 16x + 5y = 120 のグラフ。格子点（整数解）条件により唯一の解 (5, 8) が抽出される。}
\end{figure}

\clearpage
\chapter[15パズルと置換群]{第8講：15パズルの不可能性定理 —— 置換群の偶奇性と現代暗号}
\vspace{-1.0em}
\begin{center}
{\large\color{accentcrimson}\textbf{サム・ロイドの挑戦状を粉砕した対称群（Symmetric Group）の理論}}
\end{center}
\vspace{1.0em}

1870年代、アメリカ全土からヨーロッパに至るまで社会現象を巻き起こした「15パズル（15 Puzzle, \#234）」。1から15までの数字が彫られた正方形の駒を、$4 \times 4$ の枠の中で1マスの空白を利用してスライドさせ、順番通りに並べ替えるスライディング・ブロック・パズルです。

稀代のパズル考案者サム・ロイドは、「14」と「15」の駒だけが逆になった初期配置を提示し、「これを元の1から15の昇順に直すことができたら1,000ドルの賞金（現在の価値で数百万円）を支払う」という懸賞金をかけて世界中を狂乱させました。しかし、賞金を手にした者は誰一人としていませんでした。なぜなら、このパズルを解くことは**「数学的に不可能」**だったからです。

現代の代数学（群論）は、この謎を鮮やかに解き明かします。
15個の駒の並び替えは、数学的には**「16元集合上の置換（Permutation）」**として表されます。置換は、2つの要素を入れ替える基本操作「互換（Transposition）」の合成として表すことができ、互換の個数が偶数回であるものを「偶置換（Even Permutation）」、奇数回であるものを「奇置換（Odd Permutation）」と呼びます。

空きマスを1回動かす操作（駒の1回のスライド）は、置換としては互換1回に相当します。駒を動かして空きマスが最終的に元の右下隅（4行4列目）に戻るためには、上下の移動回数の和、左右の移動回数の和がそれぞれ偶数でなければなりません。したがって、スライドの総数は必ず**「偶数回」**となります。

すなわち、**盤面上の合法的な操作はすべて「偶置換」に限られ、配置のパリティ（偶奇性）は永遠に変化しません**。
サム・ロイドが提示した「14と15の逆転」は、偶置換から奇置換への遷移を要求するものであり、いかなる神業を用いても達成不可能な「不可能性の壁」に阻まれていたのです。このパリティ保存則は、現代の公開鍵暗号や誤り訂正符号の安全性の基盤をなす深遠な数学的真理です。

\vspace{1.2em}
\begin{tcolorbox}[colback=softblue,colframe=primaryblue,arc=2.5mm,boxrule=1.2pt,title={\textbf{重要ポイント：置換の偶奇性とパリティ保存則}}]
1. \textbf{対称群 $S_{16}$ の二分割}：すべての配置は「偶置換の世界（達成可能空間）」と「奇置換の世界（不可能性の空間）」に等分される。
2. \textbf{不変量（Invariant）の力}：空白マスの移動に伴うマンハッタン距離と転倒数（Inversion Number）の和の偶奇性は、不変に保たれる。
3. \textbf{不可能性の証明}：パズルを解くことと同じくらい、「解が存在しないことを論理的に証明する」ことは近代数学の決定的な精神である。
\end{tcolorbox}

\vspace{1.5em}
\begin{figure}[htbp]
\centering
\begin{tikzpicture}[scale=0.92, >=Stealth]
    % Left Box: Even Permutations (Solvable Universe)
    \node[draw=primaryblue, line width=1.5pt, fill=softblue!60, rounded corners=3mm, inner sep=10pt, text width=5.6cm, align=center] (EVEN) at (0, 0) {
        \textbf{\large 【偶置換の世界】}\\[6pt]
        \begin{tcolorbox}[colback=white, colframe=forestgreen, arc=2mm, boxrule=1pt, width=4.8cm, halign=center]
            \textbf{初期配置（偶数次）}
        \end{tcolorbox}
        \vspace{2pt}
        {\color{forestgreen}\textbf{$\updownarrow$ 相互移動可能（合法探索）}}\\[4pt]
        \begin{tcolorbox}[colback=white, colframe=forestgreen, arc=2mm, boxrule=1pt, width=4.8cm, halign=center]
            \textbf{目標配置（$1 \sim 15$ 完全整列）}
        \end{tcolorbox}
        \vspace{4pt}
        \small 全順列 $16! / 2$ 通り（可逆空間）
    };

    % Invariant Barrier (Dashed Red Line)
    \draw[line width=2.4pt, accentcrimson, dashed] (4.0, -2.5) -- (4.0, 2.5);

    % Right Box: Odd Permutations (Unsolvable)
    \node[draw=accentcrimson, line width=1.5pt, fill=accentcrimson!10, rounded corners=3mm, inner sep=10pt, text width=5.6cm, align=center] (ODD) at (8.0, 0) {
        \textbf{\large\color{accentcrimson} 【奇置換の世界】}\\[6pt]
        \begin{tcolorbox}[colback=white, colframe=accentcrimson, arc=2mm, boxrule=1.2pt, width=4.8cm, halign=center]
            \textbf{\color{accentcrimson} サム・ロイドの挑戦状}\\[2pt]
            \small （14 と 15 の逆転配置）
        \end{tcolorbox}
        \vspace{8pt}
        \textbf{\small\color{darkslate} 合法なスライド操作では}\\
        \textbf{\small\color{accentcrimson} 永久に到達不可能！}
    };

    % Barrier Label (Clear in center)
    \node[fill=white, draw=accentcrimson, line width=1.2pt, rounded corners=1mm, inner sep=4pt] at (4.0, 1.5) {
        \textbf{\small\color{accentcrimson} パリティの壁}
    };

    % Crossing Arrow with Forbidden Cross
    \draw[->, line width=2.2pt, darkslate] (ODD) -- (EVEN);
    \node[fill=white, inner sep=3pt] at (4.0, 0) {\Huge\color{accentcrimson}\textbf{$\times$}};
    \node[fill=white, draw=accentcrimson, line width=0.8pt, rounded corners=1mm, inner sep=3pt] at (4.0, -1.0) {
        \footnotesize\textbf{\color{accentcrimson} 遷移不可能}
    };

    \node[below=20pt, align=center, text width=12cm] at (4.0, -2.8) {
        \textbf{結論}：スライド操作はパリティ（置換の偶奇性）を完全に保存するため、\\2つの世界を行き来することは数学的に絶対に不可能である。
    };
\end{tikzpicture}
\caption{15パズルの状態空間分割。置換群の偶奇性（パリティ）が創り出す「超えられない不可能性の壁」。}
\end{figure}

\clearpage

% --- Epilogue ---
\chapter*{結びにかえて —— 遊びの中に宿る数学の精神}
\addcontentsline{toc}{chapter}{結びにかえて —— 遊びの中に宿る数学の精神}

ヘンリー・アーネスト・デュードニーが残したパズルの世界を巡る旅はいかがでしたでしょうか。

正三角形がヒンジで回転して正方形へと姿を変えるハバダッシャー・パズル、グラフ理論の礎となったオイラーの一筆書き、チェス盤の白黒が織りなすパリティの絶対則、そして15パズルが証明した「解けないことの美学」。

これらのパズルに共通しているのは、**「一見すると自由奔放に見える問題の中に、厳然たる数学的不変量（Invariant）が横たわっている」**という真実です。

一筆書きにおける奇頂点の個数、ナイトの移動に伴うマスの色の反転、スライドパズルにおける置換の偶奇性——どれも、どんなに試行錯誤を繰り返しても覆すことのできない「宇宙の法則」です。その法則を発見した瞬間、迷宮の霧は晴れ、混沌とした問題は極めてエレガントな秩序へと結晶化します。

大人の学び直しの喜びとは、単に公式を当てはめて計算の正解を出すことではなく、このように**「世界を貫く論理の美しさに息を呑むこと」**に他なりません。デュードニーが100年以上前に読者に届けようとしたのも、まさにその知的歓喜（Amusements）でした。

本書が、読者の皆様の知的探求を刺激し、数学という壮大で美しい知的遺産を再び楽しむきっかけとなれば幸いです。

\vspace{2.5cm}
\begin{flushright}
2026年 錦秋\\
教育情報科学研究会\\
監修・執筆陣一同
\end{flushright}

\end{document}
